\section{Related Work}
\label{sec:related-work}

\subsection{Self-Training}
\label{subsec:st}

Despite its popularity, self-training does not always improve
predictive performance, which can depend on initialization,
hyperparameters, and the unlabeled data
\cite{oliver2018realistic,su2021realistic}.
While self-training has been studied extensively in empirical work
\cite{amini2025survey,vanengelen2020survey,yang2023survey,kage2026review},
its theoretical understanding remains limited.
To address the challenges discussed in Section~\ref{sec:introduction},
existing theoretical analyses have considered various
simplifications of the self-training procedure.

For instance, Wei et al.~\cite{wei2021theoretical} establish accuracy guarantees
for fixed pseudo-labelers with consistency regularization.
Oymak and Gulcu~\cite{oymak2020} analyze confidence-thresholded
averaging estimators, using fresh batches for iterative characterization.
Frei et al.~\cite{frei2022selftraining} establish accuracy guarantees
for gradient-based self-training with fresh batches.
Other theoretical analyses investigate iterative
self-training under different structural assumptions.
Chen et al.~\cite{chen2020selftraining} study gradient-based
self-training under domain shift.
He et al.~\cite{he2022information} derive information-theoretic
generalization expressions.
Zhang et al.~\cite{zhang2022unlabeled} establish convergence and
generalization guarantees for neural-network regression.
Li et al.~\cite{li2025fixmatch} analyze FixMatch-like
gradient updates with adaptive pseudo-labels on fixed datasets.
Heidarizadeh et al.~\cite{heidarizadeh2026confidence} derive
generalization bounds under gradual domain shift.
Closer to our proportional-asymptotic setting,
Takahashi~\cite{takahashi2026role} derives sharp predictions
for linear classifiers on Gaussian mixtures, using fresh batches
and stagewise minimization.
Takanami et al.~\cite{takanami2025effect} analyze repeated-data
self-distillation with temperature-controlled soft pseudo-labels,
including hard pseudo-labels as a limiting case, and stagewise
minimization, but without confidence selection.
Both works rely on non-rigorous replica calculations from physics.

Our analysis combines gradient updates, confidence selection,
and repeated data reuse within a single rigorous framework.
The algorithm incorporates pseudo-labels refreshed at each
gradient step~\cite{lee2013pseudolabel}, hard pseudo-labeling
and confidence-based selection~\cite{sohn2020fixmatch},
and allows asymmetric confidence thresholds for the two
predicted classes, motivated by class-dependent selection
strategies~\cite{zhang2021flexmatch}.
To the best of our knowledge, no previous analysis establishes
rigorous, asymptotically exact empirical trajectory laws for
both parameter coordinates and training scores with this
combination of features in the proportional regime.

\subsection{Self-Distillation}
\label{subsec:sd}

A related procedure is self-distillation, where a student learns
from the predictions of a previously trained teacher, possibly
alongside the original labels.
While the two procedures are closely related, their objectives
typically differ.
Self-distillation aims to improve an existing estimator through
regularization, denoising, or feature transfer, whereas
semi-supervised self-training seeks to exploit the structure
of unlabeled data to improve a classifier trained on limited
labeled observations.

Self-distillation has a substantial theoretical literature,
particularly in regression, where explicit expressions or
tractable characterizations of successive minimizers allow
precise analysis through spectral and random matrix methods.
Mobahi et al.~\cite{mobahi2020selfdistillation} characterize
regularization through spectral filtering.
Das and Sanghavi~\cite{das2023understanding} study the effects
of regularization and label noise, while Pareek et
al.~\cite{pareek2024understanding} establish improvements
from repeated distillation.
Wu et al.~\cite{wu2026selftraininghelpshurtsdenoising} derive
high-dimensional risk recursions revealing a trade-off between
denoising and signal forgetting.
Further precise results are obtained by Dang et al.~\cite{dang2026}
and Lecoiu et al.~\cite{lecoiu2026}.
Classification has also been studied.
Das and Sanghavi~\cite{das2023understanding} establish accuracy
improvements for logistic classifiers under restrictive feature
assumptions, while Takanami et al.~\cite{takanami2025effect}
derive replica predictions for repeated-data Gaussian-mixture
classification.
Javanmard et al.~\cite{javanmard2025selfboost} derive 
AMP-based characterizations of iterative retraining with reused data, 
combining model predictions with noisy observed labels in 
Gaussian-mixture and generalized-linear models.

These works primarily analyze successive fitted predictors,
often exploiting special structure to obtain explicit statistical
characterizations.
Our setting instead studies the optimization dynamics themselves,
where confidence selection of hard pseudo-labels creates
discontinuities at the selection thresholds, and repeated
data reuse couples successive gradient updates.
These features prevent the direct spectral reductions available
in much of the self-distillation literature.

\subsection{Adaptive Training Procedures}
\label{subsec:adapt}

Our work relates to optimization procedures in which the
training objective depends on the current model.
In performative prediction, model deployment changes the
distribution of subsequent observations.
Perdomo et al.~\cite{perdomo2020performative} and
Mendler-D\"unner et al.~\cite{mendlerdunner2020stochastic}
establish convergence guarantees for repeated risk
minimization and stochastic optimization, with further
results on decision-dependent distributions
in~\cite{drusvyatskiy2023stochastic,cutler2024stochastic}.
Related work on adaptive sample selection includes
iterative trimmed-loss minimization, for which
Shen and Sanghavi~\cite{shen2019trimmed} establish
convergence and statistical guarantees, and
Aravkin and Davis~\cite{aravkin2020trimmed} derive
optimization complexity guarantees.
Nonconvex optimization also arises in semi-supervised
support vector machines, where several methods have been
developed~\cite{chapelle2008optimization,
collobert2006large, piccialli2024}.

Our finite-dimensional convergence analysis is related
to these approaches to adaptive and nonconvex optimization, but
we derive empirical trajectory laws for parameter coordinates
and training scores in the proportional high-dimensional regime,
accounting for adaptive hard pseudo-labels, confidence
selection, and repeated data reuse.


\subsection{High-Dimensional Dynamics}
\label{subsec:high-dim-dynamics}

The high-dimensional behavior of first-order optimization
algorithms has been studied through several approaches.
Paquette et al.~\cite{paquette2022} derive risk trajectories
for stochastic gradient descent on high-dimensional quadratic
objectives, while Ben Arous et
al.~\cite{arous2021online,arous2023highdimensional}
study scaling limits and effective dynamics of online
stochastic gradient descent in nonconvex inference problems.
Another approach originates from statistical physics, where
mean-field methods simplify systems with many interacting
degrees of freedom by describing their collective behavior
through macroscopic observables.
This viewpoint has motivated a substantial literature on
statistical inference and learning
\cite{mezard2009information,zdeborova2016statistical}.
Dynamical mean-field theory (DMFT) extends this approach
to evolving systems, describing limiting trajectories
through representative \emph{effective stochastic processes}.
When the same random data are repeatedly reused, these
processes generally involve memory terms and temporally
correlated Gaussian fluctuations.

Rigorous DMFT descriptions of first-order methods have
been established by Celentano, Cheng, and
Montanari~\cite{celentano2026highdimensionalasymptoticsordermethods}
through a reduction to approximate message passing (AMP),
and by Gerbelot et al.~\cite{gerbelot2024rigorous}
through iterative Gaussian conditioning.
Mignacco et al.~\cite{Mignacco_2021} derive dynamical
mean-field equations for gradient-based classification
of Gaussian mixtures, providing a related statistical-physics
analysis without confidence selection.
AMP is closely related to DMFT.
Its key feature is the Onsager correction, which compensates
for correlations generated by repeated matrix use,
leading to a simpler effective description through
\emph{state evolution}: deterministic recursions tracking
the limiting distributions of the iterates.
For Gaussian matrices, rigorous AMP analyses use Gaussian
conditioning, as in Bolthausen~\cite{bolthausen2009}
and Bayati and Montanari~\cite{bayati2011dynamics};
see also~\cite{feng2022unifying}.
Ranasinghe et al.~\cite{ranasinghe2026state} extend scalar
AMP state evolution to certain discontinuous denoisers, although applying
their results to our recursion requires an AMP reduction.
Yang's tensor programs provide another framework for
characterizing limiting coordinate laws of random-matrix
computations, including neural-network training
\cite{yang2020neuralmatrixlaws,yang2022}.
Applying these frameworks to our recursion requires accommodating 
hard confidence selection, evolving scalar parameters such as the 
intercept that shift the selection boundaries, 
and a mean-shifted Gaussian design.
Beyond finite-horizon characterization, the long-time
behavior of effective processes presents additional challenges.
Celentano et
al.~\cite{celentano2026highdimensionalasymptoticsordermethods}
derive fixed-point equations under exponential convergence
conditions, while Chen and
Shen~\cite{chen2025learningsingleindexmodel}
establish long-time DMFT convergence under suitable local
curvature assumptions.
Fan et al.~\cite{fan2025dynamicalmeanfieldanalysisadaptivepropagation,
fan2025dynamicalmeanfieldanalysisadaptivereplica}
study finite-horizon and equilibrium behavior of adaptive
Langevin dynamics under additional assumptions.

These frameworks provide general tools for analyzing
high-dimensional dynamics.
For our recursion, a direct application requires additional
arguments for discontinuous confidence selection combined
with evolving scalar shifts and data reuse.
We develop a self-contained Gaussian conditioning
argument tailored to our recursion, constructing the Gaussian
innovations, projections, and limiting effective processes.
We further establish long-time convergence of these processes.
